\uuid{cO4m}
\titre{Probabilités, événements}
\auteur{worms}
\organisation{exo7}
\licence{CC BY-NC-SA 4.0}
\datecreate{2017-08-28}
\niveau{L2}
\module{Probabilité et statistique}
\chapitre{Probabilité discrète}
\sousChapitre{Probabilité et dénombrement}
\difficulte{}
\theme{Probabilités, événements}
\exo7qcmid{worms.worms.020}

\contenu{
\begin{question}
\qclassification{260.01}
Soit $n$ un entier non nul et $k$ un entier compris entre $1$ et $n$. On lance $n$ fois une pièce équilibrée. Quelle est la probabilité d'obtenir exactement $k$ piles ?
\begin{answers}
\bad{$C_n^k p^n(1-p)^{n-k}$}
\bad{$\displaystyle\frac{k!}{n!(n-k)!}(1/2)^k(1/2)^{n-k}$}
\good{$n!/(2^n \, k!(n-k)!)$}
\end{answers}
\begin{explanations}
Ce n'est pas $C_n^k p^n(1-p)^{n-k}$, mais $C_n^k p^k(1-p)^{n-k}$ et comme $p=1/2$ alors
$p^k(1-p)^{n-k}= \frac{1}{2^n}$.
\end{explanations}
\end{question}
}
